\section{Problem formulation and design objectives}

An HLS is not only a portfolio of heterogeneous components to be exploited at the current time. It is an evolving organization of usable competences. Let the provisional competence portfolio be $\Ct=[c_{iz}(t)]$, where $c_{iz}(t)$ denotes effective competence of member $i$ for task region $z$. This is conceptual notation: competence need not be directly observable, and it is not assumed to be identical to accuracy or a sufficient state representation.

At time $t$, the environment $E_t$ includes the task demand and may also affect costs, latency requirements, resource availability, operational constraints, and required competences. It may be stationary or change over time. The design objective is not an artificial objective of maximizing diversity or every $c_{iz}$ independently. It is to understand when the organization of competences has functional value for long-horizon system operation.

Two linked allocation problems follow:
\begin{equation}
\text{who should solve now?}\qquad \neq \qquad \text{who should learn?}
\label{eq:two-allocations}
\end{equation}
The first selects a current division of labour. The second decides whether knowledge generated by operation or by a capable intervention should be retained, transferred, or used to transform one or more members. The member that fails, serves a task, or produces knowledge need not be the member that should learn.

The repository-level research question remains:
\begin{quote}
Can the deliberate evolution of competence allocation in a heterogeneous learning system improve its long-term performance compared with independently optimizing task allocation and knowledge transfer?
\end{quote}
It is a scientific hypothesis, not a result of this paper. The long-horizon comparison must be conditioned on a specified environment, horizon, information set, resources, and strong alternatives.

The design objectives of this draft are therefore modest: make the architecture and its candidate decisions explicit; identify importable foundations and their structural limits; preserve preliminary minimal-model evidence; and state how a future validation programme could test the general question without presuming a positive outcome.
